\documentclass[12pt, a4paper]{ctexart}
\usepackage{fontspec}
\usepackage{xcolor}
\usepackage{hyperref}
\usepackage{geometry}
\usepackage{enumitem}
\usepackage{parskip}
\usepackage{titlesec}
\usepackage{tcolorbox}
\tcbuselibrary{breakable}
\tcbset{autoparskip}

\usepackage{setspace}
\usepackage{listings}
\usepackage{amssymb}

\geometry{left=2.5cm, right=2.5cm, top=2.5cm, bottom=2.5cm}

\setmainfont{Times New Roman}
\setCJKmainfont{SimSun}

\hypersetup{
    colorlinks=true,
    linkcolor=blue!70!black,
    urlcolor=cyan,
    pdftitle={面向对象（上）-逐题精讲解义},
    pdfauthor={专升本-fifine老师}
}

\definecolor{myblue}{RGB}{30,100,200}
\definecolor{lightblue}{RGB}{230,240,255}
\definecolor{darkblue}{RGB}{0,50,120}

\newtcolorbox{bluebox}[1][]{
    breakable,
    colback=lightblue,
    colframe=myblue,
    coltitle=darkblue,
    fonttitle=\bfseries,
    title={#1},
    boxrule=0.8pt,
    arc=4pt,
    left=6pt,
    right=6pt,
    top=6pt,
    bottom=6pt,
    before skip=8pt,
    after skip=8pt
}

\usepackage{etoolbox}
\BeforeBeginEnvironment{bluebox}{\par\vfill}%
\AfterEndEnvironment{bluebox}{\par\vfill}%

\titleformat{\section}
  {\normalfont\Large\bfseries\color{myblue}}{\thesection}{1em}{}
\titlespacing*{\section}{0pt}{3.5ex plus 1ex minus .2ex}{1.5ex plus .2ex}

\title{\textcolor{myblue}{\textbf{面向对象（上） - 逐题精讲解义}}}
\author{专升本-fifine老师 教学组}
\date{2026年03月05日}

\lstset{
    language=Java,
    basicstyle=\ttfamily\small,
    keywordstyle=\color{blue},
    commentstyle=\color{green!60!black},
    stringstyle=\color{orange},
    numbers=left,
    numberstyle=\tiny\color{gray},
    frame=single,
    breaklines=true,
    columns=fullflexible,
    showstringspaces=false
}

\newcommand{\code}[1]{\texttt{#1}}

\begin{document}

\maketitle
\thispagestyle{empty}
\newpage

\section{面向对象（上）}

\begin{enumerate}[label=\textbf{\arabic*.}]

	\item \textbf{以下关于类和对象的说法正确的是}
	      \begin{enumerate}[label=\Alph*.]
		      \item 类是对象的实例
		      \item 对象是类的模板
		      \item 一个类只能创建一个对象
		      \item 对象是类的具体实例
	      \end{enumerate}

	      \begin{bluebox}{答案与深度解析}
		      \textbf{答案：D. 对象是类的具体实例}

		      \textbf{【核心认知框架】}
		      \begin{itemize}[label={}, leftmargin=*, nosep]
			      \item \textbf{题目本质}：考查类与对象的本质关系——抽象模板与具体实例
			      \item \textbf{关键区分点}：类=定义（图纸），对象=实例（楼房）
			      \item \textbf{命题人意图}：测试对面向对象核心概念的理解，反向设置混淆项
		      \end{itemize}

		      \textbf{【选项原理深度拆解】}

		      \begin{itemize}[leftmargin=*, nosep, label={}]
			      \item[A] \textbf{类是对象的实例 — 错误（概念颠倒）}
			            \begin{itemize}[label={}, leftmargin=*, nosep]
				            \item \textbf{字节码层面}：Class对象存储在方法区，不是实例
				            \item \textbf{类加载机制}：JVM加载.class文件时创建Class对象
				            \item \textbf{空间布局}：类存储在方法区，对象存储在堆内存
				            \item \textbf{字节码铁证}：
				                  \begin{lstlisting}[language=Java]
// javap -c 输出
class User {
  // 字段、方法定义在方法区的Class对象中
}
// 对象存储在堆中
User user = new User(); // user引用指向堆内存
				            \end{lstlisting}
			            \end{itemize}

			      \item[B] \textbf{对象是类的模板 — 错误（概念颠倒）}
			            \begin{itemize}[label={}, leftmargin=*, nosep]
				            \item \textbf{JVM规范}：Class对象是对象的模板（JLS §12.1）
				            \item \textbf{实例化过程}：new指令根据Class模板在堆中分配内存
				            \item \textbf{代码示例}：
				                  \begin{lstlisting}[language=Java]
// 类=模板（图纸）
class House {
    int rooms;
    String address;

    public House(int r, String a) {
        this.rooms = r;
        this.address = a;
    }
}

// 对象=实例（根据图纸建造的楼房）
House house1 = new House(3, "北京路1号");  // 第一个对象
House house2 = new House(2, "上海路2号");  // 第二个对象
// house1和house2是根据House类创建的不同实例
				            \end{lstlisting}
			            \end{itemize}

			      \item[C] \textbf{一个类只能创建一个对象 — 错误（表述不准确）}
			            \begin{itemize}[label={}, leftmargin=*, nosep]
				            \item \textbf{单例模式}：仅当类设计为单例时才只创建一个对象
				            \item \textbf{常规情况}：类可以创建多个独立对象
				            \item \textbf{堆内存验证}：
				                  \begin{lstlisting}[language=Java]
User user1 = new User();
User user2 = new User();
User user3 = new User();
// user1, user2, user3 指向堆中三个不同的对象
// 每个对象有独立的实例变量
System.out.println(user1 == user2);  // false
				            \end{lstlisting}
				            \item \textbf{命题陷阱}：将单例模式的特殊性质泛化为所有类
			            \end{itemize}

			      \item[D] \textbf{对象是类的具体实例 — 正确}
			            \begin{itemize}[label={}, leftmargin=*, nosep]
				            \item \textbf{JVM规范}：对象是类的实例化（JLS §12.5）
				            \item \textbf{实例化过程}：
				                  \begin{enumerate}[label={}, nosep]
					                  \item 类加载（Class对象创建）
					                  \item 分配堆内存 + 调用<init>方法初始化
					                  \item 返回对象引用
				                  \end{enumerate}
				            \item \textbf{字节码证据}：
				                  \begin{lstlisting}[language=Java]
// 源码
User user = new User();

// 字节码（javap -c）
0:  new           #2                  // class User
3:  dup
4:  invokespecial #3                  // Method User."<init>":()V
7:  astore_1
// new指令：分配堆内存
// invokespecial：调用构造方法初始化
// astore_1：将引用存入局部变量表
				            \end{lstlisting}
				            \item \textbf{完整代码示例}：
				                  \begin{lstlisting}[language=Java]
// 类定义（模板）
class Person {
    private String name;
    private int age;

    public Person(String name, int age) {
        this.name = name;
        this.age = age;
    }

    public void introduce() {
        System.out.println("姓名: " + name + ", 年龄: " + age);
    }
}

// 使用类创建对象（实例化）
public class Demo {
    public static void main(String[] args) {
        // 对象是Person类的具体实例
        Person p1 = new Person("张三", 25);  // 实例1
        Person p2 = new Person("李四", 30);  // 实例2

        // 每个对象是独立的存在
        p1.introduce();  // 姓名: 张三, 年龄: 25
        p2.introduce();  // 姓名: 李四, 年龄: 30
    }
}
				            \end{lstlisting}
			            \end{itemize}
		      \end{itemize}

		      \textbf{【应背诵知识点】}
		      \begin{itemize}[label={}, nosep]
			      \item 类与对象关系：类=抽象模板，对象=具体实例
			      \item 类可以创建多个对象，每个对象独立存在
			      \item 选项辨析：
			            \begin{itemize}[label={}, nosep]
				            \item "类是对象的实例"→错误，关系颠倒
				            \item "对象是类的模板"→错误，关系颠倒
				            \item "一个类只能创建一个对象"→错误（单例除外）
				            \item "对象是类的具体实例"→正确
			            \end{itemize}
		      \end{itemize}

		      \textbf{【命题趋势与实战策略】}

		      \textbf{考场秒杀三步法}：
		      \begin{enumerate}[label={}, nosep]
			      \item 看到"类"→联想"模板、图纸、抽象"
			      \item 看到"对象"→联想"实例、实例化、具体"
			      \item 类与对象关系：类创建对象（类→对象）
		      \end{enumerate}

		      \textbf{记忆口诀}：
		      \begin{itemize}[label={}, nosep]
			      \item 类是蓝图对象楼，一图多楼不固定
			      \item 实例化后堆中留，属性独立各不同
		      \end{itemize}

		      \textbf{高频陷阱}：
		      \begin{enumerate}[label={}, nosep]
			      \item 将"模板"与"实例"的关系写反
			      \item 将单例模式的特殊性泛化
			      \item 混淆类加载与对象实例化
		      \end{enumerate}

		      \textbf{5秒决策技巧}：看到"对象是类的实例"→直接选，这是面向对象的定义
	      \end{bluebox}

	      \vspace{1em}

	\item \textbf{以下哪个是Java中的对象引用传递？}
	      \begin{enumerate}[label=\Alph*.]
		      \item 基本数据类型参数传递
		      \item 对象作为参数传递
		      \item 数组作为参数传递
		      \item 以上都是
	      \end{enumerate}

	      \begin{bluebox}{答案与深度解析}
		      \textbf{答案：B. 对象作为参数传递}

		      \textbf{【核心认知框架】}
		      \begin{itemize}[label={}, leftmargin=*, nosep]
			      \item \textbf{题目本质}：考查Java参数传递的值传递本质 + 引用传递表现
			      \item \textbf{关键区分点}：基本类型=值复制，对象=引用地址复制
			      \item \textbf{命题人意图}：测试"Java是值传递"这一核心概念
		      \end{itemize}

		      \textbf{【选项原理深度拆解】}

		      \begin{itemize}[leftmargin=*, nosep, label={}]
			      \item[A] \textbf{基本数据类型参数传递 — 错误（值传递）}
			            \begin{itemize}[label={}, leftmargin=*, nosep]
				            \item \textbf{传递机制}：复制值的副本到方法栈帧
				            \item \textbf{字节码验证}：
				                  \begin{lstlisting}[language=Java]
public static void swap(int a, int b) {
    int temp = a;  // a是副本
    a = b;         // 修改副本
    b = temp;
}
// 调用
int x = 5, y = 10;
swap(x, y);        // x, y的副本参与交换
System.out.println(x); // 输出5（原值不变）
System.out.println(y); // 输出10（原值不变）
				            \end{lstlisting}
				            \item \textbf{内存模型}：
				                  \begin{itemize}[label={}, nosep]
					                  \item main方法栈帧：存储x=5, y=10
					                  \item swap方法栈帧：存储a=5（副本）, b=10（副本）
					                  \item 修改a,b不影响x,y（不同栈帧独立）
				                  \end{itemize}
				            \item \textbf{命题陷阱}：误认为是"引用传递"（Java不存在C++的引用传递）
			            \end{itemize}

			      \item[B] \textbf{对象作为参数传递 — 正确（表现类似引用传递）}
			            \begin{itemize}[label={}, leftmargin=*, nosep]
				            \item \textbf{传递机制}：传递引用地址的副本（值传递的特殊情况）
				            \item \textbf{核心概念}：传递的是对象的引用（地址），不是对象本身
				            \item \textbf{字节码证据}：
				                  \begin{lstlisting}[language=Java]
class User {
    int age;
}

public static void modify(User u) {
    u.age = 30;  // 通过引用修改对象属性
}

// 调用
User user = new User();
user.age = 20;
modify(user);   // 传递引用地址副本
System.out.println(user.age); // 输出30（被修改）
				            \end{lstlisting}
				            \item \textbf{内存模型验证}：
				                  \begin{enumerate}[label={}, nosep]
					                  \item 堆内存：User对象@0x1234，age=20
					                  \item main栈帧：user引用 → 0x1234
					                  \item modify栈帧：u引用 → 0x1234（地址副本）
					                  \item u.age直接访问堆中对象，修改成功
				                  \end{enumerate}
				            \item \textbf{关键特性验证}：
				                  \begin{lstlisting}[language=Java]
class User {
    String name;
}

// 证明是值传递（引用地址为值）
public static void reassign(User u, User newUser) {
    u = newUser;  // 修改引用副本，不影响原始引用
}

User u1 = new User();
User u2 = new User();
reassign(u1, u2);
// u1仍指向原对象（只是副本被修改）
				            \end{lstlisting}
				            \item \textbf{完整代码示例}：
				                  \begin{lstlisting}[language=Java]
class Person {
    String name;

    public Person(String name) {
        this.name = name;
    }
}

public class ReferencePassDemo {
    // 对象作为参数传递（引用地址传递）
    public static void changeName(Person p, String newName) {
        // p是引用地址的副本
        // 通过引用访问堆中的对象
        p.name = newName;  // 修改原始对象的属性
    }

    // 重新赋值引用副本
    public static void changeReference(Person p) {
        p = new Person("新对象");  // 修改副本引用，不影响原引用
    }

    public static void main(String[] args) {
        Person person = new Person("张三");
        System.out.println("改变前: " + person.name);  // 张三

        changeName(person, "李四");
        System.out.println("改变属性后: " + person.name);  // 李四

        Person original = person;
        changeReference(person);
        System.out.println("改变引用后: " + person.name);  // 李四（原引用不变）
        System.out.println("same object: " + (person == original));  // true
    }
}
				            \end{lstlisting}
			            \end{itemize}

			      \item[C] \textbf{数组作为参数传递 — 错误（与B类似）}
			            \begin{itemize}[label={}, leftmargin=*, nosep]
				            \item \textbf{传递机制}：数组是对象，传递引用地址副本
				            \item \textbf{与对象相同}：本质相同，都是对象引用传递
				            \item \textbf{代码示例}：
				                  \begin{lstlisting}[language=Java]
public static void modifyArray(int[] arr) {
    arr[0] = 100;  // 修改原数组
}

int[] array = {1, 2, 3};
modifyArray(array);
System.out.println(array[0]);  // 输出100
				            \end{lstlisting}
				            \item \textbf{命题陷阱}：D项"以上都是"可能看似正确，但A是基本类型值传递
			            \end{itemize}

			      \item[D] \textbf{以上都是 — 错误（A是值传递）}
			            \begin{itemize}[label={}, leftmargin=*, nosep]
				            \item \textbf{错误原因}：A项"基本数据类型参数传递"是值传递，不是引用传递
				            \item \textbf{易混淆点}：
				                  \begin{itemize}[label={}, nosep]
					                  \item 基本类型：值传递（值复制）
					                  \item 对象/数组：值传递（引用地址复制）
				                  \end{itemize}
				            \item \textbf{Java只有值传递}：
				                  \begin{lstlisting}[language=Java]
// Java只有值传递（Pass-by-Value）
// 基本类型：传递值的副本
int a = 10;
void method(int x) { x = 20; }  // 不影响a

// 对象类型：传递引用地址的副本
Object obj = new Object();
void method(Object o) { o = new Object(); }  // 不影响obj
void method(Object o) { o.field = 10; }       // 影响obj（通过引用修改）
				            \end{lstlisting}
			            \end{itemize}
		      \end{itemize}

		      \textbf{【应背诵知识点】}
		      \begin{itemize}[label={}, nosep]
			      \item Java参数传递规则：只有值传递，没有引用传递
			      \item 基本类型：传值副本，修改不影响原值
			      \item 对象/数组：传引用地址副本，可修改原对象但不能换对象
			      \item 选项辨析：
			            \begin{itemize}[label={}, nosep]
				            \item 基本类型→值传递
				            \item 对象作为参数→引用传递的表现
				            \item 数组作为参数→引用传递的表现
				            \item "以上都是"→错误（A不正确）
			            \end{itemize}
		      \end{itemize}

		      \textbf{【命题趋势与实战策略】}

		      \textbf{考场秒杀三步法}：
		      \begin{enumerate}[label={}, nosep]
			      \item 看到"引用传递"→找对象/数组
			      \item 看到"基本类型"→排除（值传递）
			      \item "对象作为参数传递"是标准表述
		      \end{enumerate}

		      \textbf{记忆口诀}：
		      \begin{itemize}[label={}, nosep]
			      \item 基本类型传副本，改了也不变
			      \item 对象传地址，改属性可以，改对象不行
			      \item Java只有值传递，引用传递是现象
		      \end{itemize}

		      \textbf{高频陷阱}：
		      \begin{enumerate}[label={}, nosep]
			      \item 认为"传对象=引用传递"（本质仍是值传递）
			      \item 混淆"修改对象属性"和"重新赋值引用"
			      \item D项"以上都是"看似正确但包含错误选项
		      \end{enumerate}

		      \textbf{5秒决策技巧}：看到"对象作为参数传递"→最准确表述，直接选
	      \end{bluebox}

\end{enumerate}

\end{document}
